\scanpage{79}\mspage{10 //} 
On a donc
\[
\begin{aligned}
 u&\in\mathrm{Hom}(x,\unsure{\xi|\zeta})\subset\mathrm{Hom}(x,\xi)\amalg\mathrm{Hom}(x,\xi')\overset{\rm def}{=}E(x)\\
 u'&\in\mathrm{Hom}(x,\xi')\subset E(\unsure{a|x})\\
 \text{et}\quad u&\sim u'\quad\mathrm{mod}\ R(x).
\end{aligned}
\]
Donc \deleted{il existe} comme $u\ne u'$\inserted{(en tant qu'objets de $E(x)$)}, il existe une chaine d'objets de $E(x)$
\diagram{\begin{tikzcd}[ampersand replacement=\&,row sep=small,column sep=small]
 \&\xi\&\\a\arrow[ur,dash,"\lambda"]\arrow[dr,dash,"\lambda'"']\&\&b\arrow[ul,"\mu"']\arrow[dl,"\mu'"]\\
 \&\xi'\&\\
 x\arrow[uuur,dashed,"u"']\arrow[ur,"u'"]\&\&
\end{tikzcd}}
\[f_0=u,\ f_1,\ f_2,\ \ldots\ f_n=u'\]
$(n\geq1)$, telle que pour $0\leq i\leq n-1$, $f_i$ et $f_{i+1}$ soient reliés par une des deux relations \unsure{$(\alpha)$|$(*)$|$(1)$}\inserted{ou leurs sym.}. En particulier \deleted{\unsure{on aura}} prenant $i=0$, il existe $v:x\to a$ tel que $\underset{\substack{\Vert\\u}}{f_0}=\lambda v$ (et \deleted{\unsure{$\varnothing$}} $f_1=\lambda'v$), ou bien $w:x\to b$ tel que $\underset{\substack{\Vert\\u}}{f_0}=\mu w$, (et $f_1=\mu'w$).

Donc $u$ se factorise soit par $(a,\lambda)$, soit par $(b,\mu)$. Mais cela signifie aussi que $\tilde u:\tilde x\mathrel{\unsure{\leftarrow|-|\to}}\tilde\xi$ se factorise soit par $\tilde a$, soit par $\tilde b$, \deleted{\unsure{qed.}} \deleted{donc que} \deleted{$\tilde\xi$ est} $\tilde\xi$ est dans $\widetilde X_{//\tilde a}\cup\widetilde X_{//\tilde b}$,
qed.
% archive sheet number: 78

